\documentclass[12pt, a4paper]{article}
\usepackage[utf8]{inputenc}
\usepackage{amsmath, amssymb, array}
\usepackage{geometry}
\geometry{a4paper, margin=1in}
\usepackage{booktabs}
\usepackage{graphicx}
\usepackage{hyperref}
\usepackage{parskip}

\title{\textbf{Sunina: The Language of Truth}}
\author{A Universal Constructed Language}
\date{}

\begin{document}

\maketitle

\textbf{Literal Translation:} Ordered-mind speech (\textit{su} = good/order, \textit{ni} = mind/focus, \textit{na} = speak) \\
\textbf{Core Philosophy:} A universal, mathematically consistent language engineered for structural clarity, profound emotional expression, and logical compositionality. Every word decomposes into familiar roots that hint at its meaning --- making vocabulary \textbf{memorable by design}, not arbitrary.

\tableofcontents
\bigskip

\section{Phonology (The Universal Sounds)}
Sunina uses only 11 sounds, chosen because they exist in nearly every human language. It strictly follows a \textbf{CV (Consonant-Vowel)} structure.
\begin{itemize}
    \item \textbf{8 Consonants:} p /p/, t /t/, k /k/, m /m/, n /n/, s /s/, w /w/, y /j/
    \item \textbf{3 Vowels:} a /a/ (\textit{father}), i /i/ (\textit{machine}), u /u/ (\textit{flute})
\end{itemize}

\subsection{Prosody}
\begin{itemize}
    \item \textbf{Rhythm:} Syllable-timed --- every syllable receives equal duration (like Spanish or Hindi).
    \item \textbf{Stress:} primary stress falls on the \textbf{first syllable} of the word, with secondary stress on the first syllable of subsequent roots within the compound. In \textit{supi}, stress is on \textbf{su}. In \textit{masupi}, primary stress is on \textbf{ma} and secondary stress on \textbf{su}. This maintains morpheme clarity without being ambiguous.
    \item \textbf{Intonation:} Statements fall in pitch at the end. Questions (containing \textit{pa}) rise on the final syllable.
\end{itemize}

\subsection{Morphophonology}
No sound changes occur at morpheme boundaries. Each root retains its canonical CV form in all positions: \textit{mu-ti-mu} keeps all three syllables distinct. This guarantees that roots are always recoverable by listeners, at the cost of longer words.

\section{The Periodic Table of Meaning (24 Roots)}
Every concept is built by combining these 24 prime syllables. The consonant defines the semantic domain; the vowel modulates it.

\vspace{0.5em}
\noindent\resizebox{\textwidth}{!}{%
\begin{tabular}{@{}llll@{}}
\toprule
\textbf{Consonant} & \textbf{-a (Open/Base)} & \textbf{-i (High/Internal)} & \textbf{-u (Round/Forward)} \\ \midrule
\textbf{P} (Life/Unknown) & \textbf{pa:} Question & \textbf{pi:} Life / Person & \textbf{pu:} The Unknown \\
\textbf{T} (Time/Target)  & \textbf{ta:} Present / Time & \textbf{ti:} Past / Behind & \textbf{tu:} You / Target \\
\textbf{K} (Action/System)& \textbf{ka:} Do / Make / Equal & \textbf{ki:} Bond / Rule / System & \textbf{ku:} Force / Harshness \\
\textbf{M} (Self/Motion)  & \textbf{ma:} Opposite / No & \textbf{mi:} I / Me / Self & \textbf{mu:} Move / Go / Change \\
\textbf{N} (Thought/Out)  & \textbf{na:} Speak / Word & \textbf{ni:} Mind / Think & \textbf{nu:} Future / Forward \\
\textbf{S} (State/Reality)& \textbf{sa:} Thing / Reality & \textbf{si:} Feel / Sense & \textbf{su:} Good / Positive \\
\textbf{W} (Experience/Magnitude) & \textbf{wa:} Intense / Big & \textbf{wi:} Weight / Sorrow & \textbf{wu:} Consume / Take in \\
\textbf{Y} (Environment/Spirit) & \textbf{ya:} Joy / Harmony & \textbf{yi:} Space / Light & \textbf{yu:} Fluid / Nature \\ \bottomrule
\end{tabular}}

\section{Grammar \& Syntax}
Sunina uses strict SOV word order, agglutination (stacking roots), and a minimal set of particles derived from the root system.

\subsection{Word Order: SOV (Subject -- Object -- Verb)}
\begin{itemize}
    \item \textit{English:} I see the friend. \quad $\rightarrow$ \quad \textit{Sunina:} \textbf{mi supi yiwu.} \quad (I friend light-consume.) \\
    \textit{Script:} $[\ominus'] \; [Z\_] \; [\bigcirc'] \; [\sqcup'] \; [\smile\_]$
\end{itemize}

\subsection{Word Building (Agglutination)}
Stack roots to build the literal essence of a concept. All compounds are \textbf{head-final}: the last root is the semantic head (what the word \textit{is}), and preceding roots modify it. This mirrors SOV typology, where the core element comes last.
\begin{itemize}
    \item \textbf{supi} = Friend (\textit{su} [good] + \textit{pi} [\textbf{person}] $\rightarrow$ ``a good person'')
    \item \textbf{masupi} = Enemy (\textit{ma} [opposite] + \textit{supi} [\textbf{friend}])
    \item \textbf{yusa} = Water (\textit{yu} [fluid] + \textit{sa} [\textbf{thing}] $\rightarrow$ ``a fluid thing'')
    \item \textbf{nisa} = Book (\textit{ni} [mind] + \textit{sa} [\textbf{thing}])
    \item \textbf{yiwu} = See (\textit{yi} [light] + \textit{wu} [\textbf{take in}] $\rightarrow$ ``take in light'')
\end{itemize}
\textit{Compound-Order Rule:} In any compound of $n$ roots, the rightmost root is the head; all roots to its left narrow its meaning, outermost modifier first. E.g., \textit{wa-ku-pi-sa} = intense $>$ force $>$ person $>$ thing $=$ ``fierce creature.''

\textit{Shortening Rule:} The official short form of a compound is always its \textbf{first two syllables} (the first stressed foot), making short forms predictable. The Sunina Academy approves short forms for any standardized compound of \textbf{3 or more syllables}; until approved, the full form is mandatory. E.g., \textit{wakupisa} (fierce creature) $\rightarrow$ \textit{waku}.

\textit{Collision Avoidance:} A short form must never duplicate any standardized lexical item, root, grammatical particle, or previously approved short form. If the first two syllables collide, the short form extends to the first three syllables instead. If that still collides, or if the three-syllable form is no shorter than the full word, the word has no short form.

\textit{Semantic Narrowing:} Compound words provide a \textbf{semantic neighbourhood}, not an exact definition. \textit{supi} (good + person) locates the concept in the zone of ``positive human relationship'' --- its precise meaning (\textit{friend}) is established by the standardized lexicon. This is by design: decomposition makes vocabulary \textbf{memorable}, not self-explanatory. Much like English \textit{understand} is not literally ``standing under,'' Sunina compounds are mnemonic scaffolds that aid learning and recall.

\textit{Fused vs.\ Free Distinction:} A fused compound (written as one word) carries its \textbf{standardized meaning}. The same roots written as separate words retain their literal, compositional sense. This lets speakers distinguish specialized intent from descriptive intent:
\begin{itemize}
    \item \textbf{supi} (one word) = friend (standardized)
    \item \textbf{su pi} (two words) = a good person (literal)
    \item \textbf{masu supi} = a bad friend (modifier + standardized compound)
\end{itemize}
Stress and writing reinforce this: fused compounds receive primary stress on the first syllable as a single unit, while free modifier phrases stress each word independently.

\subsection{Pronouns}
\begin{itemize}
    \item \textbf{mi} $[\ominus']$ = I / me \quad \textbf{tu} $[\triangle\_]$ = you \quad \textbf{pi} $[\bigcirc']$ = he / she / they \quad \textbf{sa} $[Z\cdot]$ = it
\end{itemize}
\textit{Note:} \textit{pi} and \textit{sa} naturally double as third-person pronouns from their root meanings (person, thing). Sunina makes no gender distinction.

\subsection{Plurals \& Possession}
\begin{itemize}
    \item \textbf{Plural suffix \textit{-wa}} (big $\rightarrow$ many): \quad \textit{pi} (person) $\rightarrow$ \textit{piwa} (people). \quad \textit{nisa} $\rightarrow$ \textit{nisawa} (books).
    \item \textbf{Possessive suffix \textit{-ki}} (bond) on the possessor: \quad \textit{miki supi} = my friend. \quad \textit{tuki nisa} = your book. \\
    \textit{Script (my friend):} $[\ominus'] \; [\square'] \; [Z\_] \; [\bigcirc']$
\end{itemize}
\textit{Note:} Possessive \textit{-ki} fuses with the possessor as one word (\textit{miki}, stress on \textbf{mi}). The conjunction \textit{ki} (``and'') stands alone as a separate word with its own stress.

\subsection{Grammatical Particles}
Existing roots repurposed in grammatical positions. Role is determined by sentence context.

\vspace{0.5em}
\begin{tabular}{@{}lll@{}}
\toprule
\textbf{Function} & \textbf{Particle} & \textbf{Root Logic} \\ \midrule
And & \textbf{ki} & Bond \\
But / However & \textbf{maki} & Opposite-bond \\
Or & \textbf{paki} & Question-bond \\
Because & \textbf{kiti} & Bond-past (the reason behind) \\
That / Which / Who & \textbf{ni} & Mind-link \\
When / While & \textbf{taki} & Time-bond (the linked moment) \\
If (conditional) & \textbf{panu} & Question-forward \\
\bottomrule
\end{tabular}

\vspace{0.5em}
\noindent\textit{Examples:} \\
\textit{mi ki tu supi ka.} \quad (I and you friends are.) \\
\textit{Script:} $[\ominus'] \; [\square'] \; [\triangle\_] \; [Z\_] \; [\bigcirc'] \; [\square\cdot]$ \\[0.3em]
\textit{pi ni mu-ti supi ka.} \quad (Person who came friend is.) \\
\textit{Script:} $[\bigcirc'] \; [\overline{\triangle}'] \; [\ominus\_] \; [\triangle'] \; [Z\_] \; [\bigcirc'] \; [\square\cdot]$

\subsection{Question Words}
All wh-questions begin with \textbf{pa} (question root), combined with a domain root.

\vspace{0.5em}
\begin{tabular}{@{}lll@{}}
\toprule
\textbf{English} & \textbf{Sunina} & \textbf{Logic} \\ \midrule
What? & \textbf{pasa} & question + thing \\
Who? & \textbf{papi} & question + person \\
Where? & \textbf{payi} & question + space \\
When? & \textbf{pata} & question + time \\
Why? & \textbf{pani} & question + mind \\
How? & \textbf{pamu} & question + change \\
\bottomrule
\end{tabular}

\vspace{0.5em}
\noindent\textit{Wh-example:} \textit{tu payi mu?} \quad (You where go? $=$ Where are you going?) \\
\textit{Script:} $[\triangle\_] \; [\bigcirc\cdot] \; [\sqcup'] \; [\ominus\_] \; (\bigcirc\cdot)$

\noindent\textbf{Yes/No (Polar) Questions:} Place the particle \textbf{pa} at the \textit{end} of a declarative sentence. Rising intonation on the final syllable reinforces the question.
\begin{itemize}
    \item \textit{tu mu-ti pa?} \quad (You went? $=$ Did you go?) \\
    \textit{Script:} $[\triangle\_] \; [\ominus\_] \; [\triangle'] \; [\bigcirc\cdot] \; (\bigcirc\cdot)$
    \item \textit{tasa yusa ka pa?} \quad (This water is? $=$ Is this water?)
\end{itemize}
\textit{Note:} Wh-words appear \textit{within} the clause (in the position of the unknown element), while polar \textit{pa} appears strictly clause-finally. The two patterns never overlap.

\subsection{Demonstratives}
\textbf{ta} (present $\rightarrow$ proximity) and \textbf{ti} (past/behind $\rightarrow$ distance) combine with domain roots.

\vspace{0.5em}
\begin{tabular}{@{}lll@{}}
\toprule
\textbf{English} & \textbf{Sunina} & \textbf{Logic} \\ \midrule
This (thing) & \textbf{tasa} & present + thing \\
That (thing) & \textbf{tisa} & behind + thing \\
Here & \textbf{tayi} & present + space \\
There & \textbf{tiyi} & behind + space \\
\bottomrule
\end{tabular}

\vspace{0.5em}
\noindent\textit{Example:} \textit{mi tasa wu.} \quad (I this-thing consume. $=$ I eat this.) \\
\textit{Script:} $[\ominus'] \; [\triangle\cdot] \; [Z\cdot] \; [\smile\_]$

\subsection{Tense (Verb Suffixes)}
\begin{itemize}
    \item \textbf{Present (unmarked):} \textit{mi sa ka.} \quad (I make the thing.) \\
    \textit{Script:} $[\ominus'] \; [Z\cdot] \; [\square\cdot]$
    \item \textbf{Past \textit{-ti}:} \textit{mi sa ka-ti.} \quad (I made the thing.) \\
    \textit{Script:} $[\ominus'] \; [Z\cdot] \; [\square\cdot] \; [\triangle']$
    \item \textbf{Future \textit{-nu}:} \textit{mi sa ka-nu.} \quad (I will make the thing.) \\
    \textit{Script:} $[\ominus'] \; [Z\cdot] \; [\square\cdot] \; [\overline{\triangle}\_]$
\end{itemize}
\textit{Passive:} Attach the suffix \textbf{-maka} (opposite-do) to the verb stem, before tense. The full verb template is \textbf{Verb(-Voice)-Tense(-Aspect)}. \quad \textit{sa ka-maka-ti.} $=$ The thing was made (by someone). Binding voice to the verb prevents confusion with the free negation particle \textit{ma}. \\
\textit{Copula:} \textit{ka} functions as ``is/am/are'' when linking a subject to a \textbf{noun} (classification): \\
\textit{pi supi ka.} $=$ The person is a friend. \quad \textit{Script:} $[\bigcirc'] \; [Z\_] \; [\bigcirc'] \; [\square\cdot]$ \\
\textit{Stative:} The particle \textbf{sa} (reality) before a modifier marks the subject as being \textit{in a state}, distinguishing essential quality from active performance: \\
\textit{pi sa su.} $=$ The person is good (inherent quality). \quad \textit{pi su ka.} $=$ The person does good (acts kindly).


\subsection{Aspect (Optional Verb Suffixes)}
Aspect describes \textit{how} an action unfolds over time. It is \textbf{optional} --- omit it when the distinction is unnecessary, preserving efficiency. When used, aspect stacks after tense (and optional voice): \textbf{Verb(-Voice)-Tense-Aspect}.

\vspace{0.5em}
\begin{tabular}{@{}llll@{}}
\toprule
\textbf{Aspect} & \textbf{Suffix} & \textbf{Root Logic} & \textbf{Meaning} \\ \midrule
Completed & \textbf{-ka} & do/make $=$ done & Action is finished \\
Ongoing & \textbf{-mu} & move $=$ in motion & Action is in progress \\
Habitual & \textbf{-ta} & time $=$ recurring & Action repeats regularly \\
\bottomrule
\end{tabular}

\vspace{0.5em}
\noindent\textit{Examples:} \\
\textit{mi mu-ti-ka.} \quad (I went [and arrived].) \quad
\textit{Script:} $[\ominus'] \; [\ominus\_] \; [\triangle'] \; [\square\cdot]$ \\
\textit{mi mu-ti-mu.} \quad (I was going [still on the way].) \quad
\textit{Script:} $[\ominus'] \; [\ominus\_] \; [\triangle'] \; [\ominus\_]$ \\
\textit{mi mu-ti-ta.} \quad (I used to go [regularly].) \quad
\textit{Script:} $[\ominus'] \; [\ominus\_] \; [\triangle'] \; [\triangle\cdot]$ \\
\textit{pi na-ti-mu taki yumu ka-ti.} \quad (The person was speaking when rain came.) \\
\textit{Script:} $[\bigcirc'] \; [\overline{\triangle}\cdot] \; [\triangle'] \; [\ominus\_] \; [\triangle\cdot] \; [\square'] \; [\sqcup\_] \; [\ominus\_] \; [\square\cdot] \; [\triangle']$

\vspace{0.3em}
\textit{Note:} Aspect is optional. \textit{mi mu-ti.} (I went) remains valid when the completion/duration distinction is irrelevant. Use aspect only when it adds meaning --- this keeps speech efficient.

\subsection{Affect Particles (Speaker Stance)}
Sunina expresses \textit{what} the speaker feels through modifier roots (\S3.11). To mark \textit{how} something is said --- the speaker's stance or tone --- a small set of \textbf{clause-final affect particles} is available. All are \textit{si}-initial compounds (\textit{si} = feel), making them instantly recognizable as tone markers. They appear \textit{after} the complete verb phrase (including any tense/aspect suffixes).

\vspace{0.5em}
\begin{tabular}{@{}llll@{}}
\toprule
\textbf{Stance} & \textbf{Particle} & \textbf{Logic} & \textbf{Effect} \\ \midrule
Tenderness & \textbf{siyu} & feel + nature & Said with warmth / gentleness \\
Urgency & \textbf{siku} & feel + force & Said with pressing importance \\
Regret & \textbf{sima} & feel + opposite & Said with sorrow for what happened \\
Surprise & \textbf{sipu} & feel + unknown & Said with astonishment \\
Reassurance & \textbf{siya} & feel + harmony & Said to comfort / encourage \\
\bottomrule
\end{tabular}

\vspace{0.5em}
\noindent\textit{Examples:} \\
\textit{mi mu-nu siku.} \quad (I will go --- urgently.) \\
\textit{pi ma mu-ti sima.} \quad (The person did not go --- regretfully.) \\
\textit{tu sa su siya.} \quad (You are good --- said reassuringly.) \\
\textit{mi tu yiwu-ti sipu.} \quad (I saw you --- with surprise!)

\vspace{0.3em}
\textit{Note:} Affect particles are always optional. Omitting them produces a neutral, unmarked tone. They express the speaker's attitude toward the \textit{utterance}, not the subject's emotion (which is expressed via modifier roots like \textit{ya}, \textit{wi}, \textit{ku}).

\subsection{Adjectives \& Adverbs (Free Modifiers)}
Roots can function as \textbf{free modifiers} (adjectives, adverbs, or emotions). To ensure absolute clarity, any root used as a modifier must be placed \textit{immediately before} the word it describes.

\begin{itemize}
    \item \textit{\textbf{su} pi wuka.} \quad (The good person eats.)

    \item \textit{mi \textbf{ya} mu.} \quad (I joyfully go.) \\
    \textit{Script:} $[\ominus'] \; [\sqcup\cdot] \; [\ominus\_]$
    \item \textit{mi \textbf{wi} na.} \quad (I sadly speak.) \\
    \textit{Script:} $[\ominus'] \; [\smile'] \; [\overline{\triangle}\cdot]$
    \item \textit{pi \textbf{ku} na-ti.} \quad (The person angrily spoke.) \\
    \textit{Script:} $[\bigcirc'] \; [\square\_] \; [\overline{\triangle}\cdot] \; [\triangle']$
    \item \textbf{Intensity:} Prefix \textit{wa}: \quad \textit{mi \textbf{wa wi} na.} \quad (I speak with deep sorrow.) \\
    \textit{Script:} $[\ominus'] \; [\smile\cdot] \; [\smile'] \; [\overline{\triangle}\cdot]$
    \item \textbf{Mixed emotions:} Stack freely: \quad \textit{mi \textbf{ya wi} si.} \quad (I feel bittersweet.) \\
    \textit{Script:} $[\ominus'] \; [\sqcup\cdot] \; [\smile'] \; [Z']$
    \item \textbf{Emotional intensifier:} \textit{sisi} (feel-feel) frames a clause as pure, overwhelming emotion: \\
    \textit{mi sisi ya wi si-ti.} \quad (I felt an overwhelming bittersweetness.) \\
    Use \textit{sisi} only when the speaker wishes to mark the utterance as deeply emotional, not merely descriptive.
\end{itemize}

\subsection{Comparatives}
\begin{itemize}
    \item \textbf{More:} \textit{wa} before adjective. \quad \textit{wa su} = better. \quad \textit{wa ku} = stronger.
    \item \textbf{Most:} \textit{wawa} before adjective. \quad \textit{wawa su} = best.
    \item \textbf{Less:} \textit{mawa} before adjective. \quad \textit{mawa su} = worse.
    \item \textbf{Than:} \textit{timu} (behind-move $=$ ``beyond'') after the compared noun: \quad \textit{mi tu timu wa su ka.} \quad (I, you-beyond, more good am.) \\
    \textit{Script:} $[\ominus'] \; [\triangle\_] \; [\triangle'] \; [\ominus\_] \; [\smile\cdot] \; [Z\_] \; [\square\cdot]$
\end{itemize}
\textit{Note:} \textit{ti} (past/behind) serves as a tense suffix and distance demonstrative, unified by the core metaphor of ``what is behind.'' The comparative ``than'' uses the dedicated compound \textit{timu} (beyond) to avoid overloading \textit{ti}.

\subsection{Conditionals \& Modality}
\textbf{Conditionals:} \textit{panu} (question-forward) at the start of a clause marks it as conditional.
\begin{itemize}
    \item \textit{panu yumu ka, mi mu-nu.} \quad (If rain happens, I will leave.) \\
    \textit{Script:} $[\bigcirc\cdot] \; [\overline{\triangle}\_] \; [\sqcup\_] \; [\ominus\_] \; [\square\cdot] \quad [\ominus'] \; [\ominus\_] \; [\overline{\triangle}\_]$
\end{itemize}

\noindent\textbf{Modality:} Compound particles placed before the main verb.
\begin{itemize}
    \item \textbf{suka} (good-do $=$ able): \quad \textit{mi suka mu.} \quad (I can go.) \\
    \textit{Script:} $[\ominus'] \; [Z\_] \; [\square\cdot] \; [\ominus\_]$
    \item \textbf{kuka} (force-do $=$ obliged): \quad \textit{mi kuka ka.} \quad (I must do this.) \\
    \textit{Script:} $[\ominus'] \; [\square\_] \; [\square\cdot] \; [\square\cdot]$
    \item \textbf{puka} (unknown-do $=$ uncertain): \quad \textit{mi puka mu.} \quad (I might go.) \\
    \textit{Script:} $[\ominus'] \; [\bigcirc\_] \; [\square\cdot] \; [\ominus\_]$
\end{itemize}

\subsection{Negation}
To negate an action, place the word \textbf{ma} (opposite/no) directly before the verb. This explicitly negates the clause, distinguishing it from compound antonyms (like \textit{ma-supi} = enemy).
\begin{itemize}
    \item \textit{mi ma mu.} \quad (I do not go.) \\
    \textit{Script:} $[\ominus'] \; [\ominus\cdot] \; [\ominus\_]$
\end{itemize}

\subsection{Spatial Postpositions}
Spatial relationships are marked by \textbf{postpositions} --- roots placed \textit{after} the noun they govern and before the verb, consistent with SOV order. Position in the sentence identifies their grammatical role.

\vspace{0.5em}
\begin{tabular}{@{}lll@{}}
\toprule
\textbf{English} & \textbf{Sunina} & \textbf{Root Logic} \\ \midrule
to / toward & \textbf{nu} & forward \\
from & \textbf{timu} & behind + move \\
at / in & \textbf{yi} & space \\
with (accompany) & \textbf{ki} & bond \\
\bottomrule
\end{tabular}

\vspace{0.5em}
\noindent\textit{Examples:} \\
\textit{pi yusa nu mu-ti.} \quad (Person water-toward went. $=$ The person walked to the water.) \\
\textit{mi nisa yi ni.} \quad (I book-at think. $=$ I think about the book.) \\
\textit{mi supi ki mu.} \quad (I friend-with go. $=$ I go with a friend.)

\subsection{Imperative Mood}
Commands use the particle \textbf{tuka} (target-do $=$ ``direct to act'') at the start of the sentence, preserving SOV order with zero exceptions. The subject \textit{tu} (you) is implied and may be omitted.
\begin{itemize}
    \item \textit{tuka mu!} \quad (Go!)
    \item \textit{tuka ma wuka!} \quad (Don't eat!)
    \item \textit{tuka tasa yiwu!} \quad (Look at this!)
\end{itemize}
\textit{Polite request:} Prefix \textit{su} (good): \quad \textit{su tuka na.} \quad (Please speak.)

\textit{Note:} The imperative uses \textit{tuka} (target + do), not \textit{paka}, to avoid collision with the question word \textit{pamu} (how?). Each grammar particle occupies a unique form.

\subsection{Reflexive \& Reciprocal}
\begin{itemize}
    \item \textbf{Reflexive \textit{kami}} (equal + self $=$ ``oneself''): placed in the object position. \\
    \textit{pi kami yiwu.} \quad (The person sees oneself.) \quad \textit{mi kami yiwu.} \quad (I see myself.)
    \item \textbf{Reciprocal \textit{kimi}} (bond + self $=$ ``each other''): placed in the object position. \\
    \textit{piwa kimi yiwu.} \quad (People see each other.)
\end{itemize}

\subsection{Derivational Morphology}
Systematic suffixes convert roots between grammatical categories. Because every suffix is itself a root with transparent meaning, speakers can apply them productively to any stem.

\vspace{0.5em}
\noindent\resizebox{\textwidth}{!}{%
\begin{tabular}{@{}lllll@{}}
\toprule
\textbf{Conversion} & \textbf{Suffix} & \textbf{Logic} & \textbf{Example} & \textbf{Attested} \\ \midrule
Verb $\rightarrow$ Noun (act/result) & \textbf{-sa} & thing & \textit{mu} $\rightarrow$ \textit{musa} (journey) & 30+ \\
Verb $\rightarrow$ Agent (doer) & \textbf{-pi} & person & \textit{na} $\rightarrow$ \textit{napi} (speaker) & 21 \\
Adj $\rightarrow$ Noun (quality) & \textbf{-sa} & thing & \textit{su} $\rightarrow$ \textit{susa} (goodness) & 10+ \\
Noun/Adj $\rightarrow$ Verb (become/do) & \textbf{-mu} & change & \textit{su} $\rightarrow$ \textit{sumu} (improve) & 5+ \\
Noun $\rightarrow$ Action verb & \textbf{-ka} & do & \textit{mu} $\rightarrow$ \textit{muka} (go/walk) & 30+ \\
Root $\rightarrow$ Feeling/state & \textbf{-si} & feel & \textit{ya} $\rightarrow$ \textit{yasi} (happy) & 14 \\
Any $\rightarrow$ Antonym & \textbf{ma-} & opposite & \textit{supi} $\rightarrow$ \textit{masupi} (enemy) & 30+ \\
Any $\rightarrow$ Intensified & \textbf{wa-} & intense & \textit{yusa} $\rightarrow$ \textit{wayusa} (ocean) & 10+ \\
\bottomrule
\end{tabular}}

\vspace{0.5em}
\noindent\textit{Note:} Roots used directly as pre-modifiers require no suffix (see \S3.11). The ``Attested'' column shows how many standardized words already use each pattern, confirming its productivity.

\subsection{Productive Compound Templates}
Beyond single-suffix derivation (\S3.18), Sunina's lexicon reveals recurring \textbf{multi-root templates} --- compound endings that reliably produce words in a semantic domain. These templates are now formalized as productive rules: speakers may coin new words by slotting any appropriate modifier into the \textbf{X} position without waiting for Academy approval, provided the result obeys CV phonology and passes the collision check.

\vspace{0.5em}
\noindent\textbf{Tier 1 --- Freely Productive} (high attestation, fully regular; coin at will)

\noindent\resizebox{\textwidth}{!}{%
\begin{tabular}{@{}llllp{6cm}@{}}
\toprule
\textbf{Template} & \textbf{Ending} & \textbf{Means} & \textbf{Attested} & \textbf{Examples} \\ \midrule
X-\textit{pi} & person & ``person of/who X'' & 21 & \textit{supi} (friend), \textit{nanipi} (teacher), \textit{kamupi} (worker) \\
X-\textit{ka} & do & ``to do X'' & 30+ & \textit{muka} (go), \textit{yisika} (look), \textit{sakika} (build) \\
X-\textit{sa} & thing & ``the thing of X'' & 30+ & \textit{yusa} (water), \textit{ninusa} (idea), \textit{katusa} (success) \\
X-\textit{si} & feel & ``the feeling/quality of X'' & 14 & \textit{yasi} (happy), \textit{wisi} (sad), \textit{kusi} (hard) \\
\textit{ma}-X & opposite & ``anti-/un-/non- X'' & 30+ & \textit{masupi} (enemy), \textit{mayumu} (drought) \\
\textit{wa}-X & intense & ``great/intense X'' & 10+ & \textit{wayusa} (ocean), \textit{wawisa} (mountain) \\
\bottomrule
\end{tabular}}

\vspace{0.8em}
\noindent\textbf{Tier 2 --- Domain Productive} (regular within their semantic field; coin freely for that domain)

\noindent\resizebox{\textwidth}{!}{%
\begin{tabular}{@{}llllp{6cm}@{}}
\toprule
\textbf{Template} & \textbf{Ending} & \textbf{Means} & \textbf{Attested} & \textbf{Examples} \\ \midrule
X-\textit{muka} & move+do & ``to move in X way'' & 14 & \textit{kumuka} (run), \textit{yumuka} (swim), \textit{yinumuka} (climb) \\
X-\textit{kusa} & force+thing & ``tool/instrument for X'' & 8 & \textit{wakusa} (weapon), \textit{yikakusa} (key), \textit{susikusa} (broom) \\
X-\textit{pisa} & life+thing & ``body part of X function'' & 18 & \textit{kapisa} (hand), \textit{nipisa} (head), \textit{sipisa} (heart) \\
X-\textit{pimu} & life+move & ``creature/animal of X'' & 8 & \textit{yimupimu} (bird), \textit{yumupimu} (fish), \textit{sukipimu} (dog) \\
X-\textit{wusa} & consume+thing & ``food/edible of X'' & 7 & \textit{yuwusa} (fruit), \textit{wisawusa} (grain) \\
X-\textit{yusa} & fluid+thing & ``liquid/body-of-water of X'' & 10 & \textit{muyusa} (river), \textit{wayusa} (ocean), \textit{puyusa} (fog) \\
X-\textit{kisa} & bond+thing & ``structure/building of X'' & 6 & \textit{yikisa} (house), \textit{yinukisa} (roof) \\
X-\textit{sasa} & thing+thing & ``material/substance of X'' & 4 & \textit{kuwisasa} (metal), \textit{nisasa} (paper) \\
\bottomrule
\end{tabular}}

\vspace{0.8em}
\noindent\textbf{How to read the templates.} Each template is head-final: the ending defines \textit{what the word is} (person, tool, creature\ldots), and the X-slot narrows its meaning. Multiple modifier roots may stack before the ending: \textit{wa-ki-pi} = intense $>$ rule $>$ \textbf{person} = ``ruler.''

\vspace{0.3em}
\noindent\textit{Coining example:} To create ``singer'' (a person who sings), apply the \textbf{agent} template to the verb ``sing'' (\textit{yanaka}): drop the verb ending and attach \textit{-pi} $\rightarrow$ \textit{yanapi} (sing-person). Or compose directly: \textit{ya-na-ka-pi} (joy + speak + do + person). Both routes converge on a suggestive, head-final compound.

\vspace{0.3em}
\noindent\textbf{Coining Checklist} (apply before using any newly coined word):
\begin{enumerate}
    \item \textbf{CV compliance} --- every syllable is one consonant + one vowel.
    \item \textbf{Head-final} --- the semantic head (template ending) is the rightmost element.
    \item \textbf{Collision-free} --- the full form and its short form (first 2 syllables) do not duplicate any existing word, root, particle, or approved short form. If the 2-syllable short collides, extend to 3 syllables; if that also collides, the word has no short form.
    \item \textbf{Suggestive} --- a listener who knows the 24 roots can infer the \textit{semantic neighbourhood} of the compound. The decomposition should hint at the meaning, even if the precise sense requires learning the standardized definition.
    \item \textbf{Minimal} --- use the fewest roots that uniquely identify the concept; do not over-specify.
\end{enumerate}
Words coined via Tier~1 templates are immediately usable. Words coined via Tier~2 templates are usable within their domain. Wholly novel template endings (not listed above) require Academy standardization before adoption.

\subsection{Disambiguation Conventions}
Several roots serve multiple grammatical roles. Sunina resolves potential ambiguity through strict writing and prosodic conventions:
\begin{enumerate}
    \item \textbf{Suffixes are hyphenated} to their host word: \textit{ka-ti} (made), \textit{mu-nu} (will go). They share the host word's stress domain.
    \item \textbf{Compound words are fused} into one unhyphenated word: \textit{supi} (friend), \textit{timu} (from/than), \textit{kuka} (must).
    \item \textbf{Free particles stand alone} as separate words with independent stress: \textit{ki} (and), \textit{ma} (not), \textit{panu} (if).
\end{enumerate}
Thus negation \textit{ma ka-ti} (not $+$ did) is visually and prosodically distinct from the passive suffix \textit{ka-maka-ti} (bound to the verb within its stress domain). Similarly, the conjunction \textit{ki} (standalone) is distinct from the possessive suffix \textit{-ki} (attached to the possessor).

\newpage

\section{Core Vocabulary}
The following words are \textbf{standardized compounds} derived from the 24 roots. They are the official forms --- speakers should prefer these over ad-hoc coinages to ensure mutual intelligibility. Short forms (where listed) may be used in casual speech.

\subsection{People \& Relations}
\begin{tabular}{@{}llll@{}}
\toprule
\textbf{English} & \textbf{Sunina} & \textbf{Short} & \textbf{Logic} \\ \midrule
Friend & supi & --- & good + person \\
Enemy & masupi & --- & opposite + friend \\
Child & nupi & --- & future + person \\
Elder & tipi & --- & past + person \\
Baby / Infant & mawanupi & mawanu & opposite-intense + child (small child) \\
Parent & kipi & --- & bond + person \\
Family & kipiwa & --- & bond + people \\
Partner / Spouse & kisupi & kisu & bond + good + person \\
Sibling & kakipi & kaki & equal-bond + person \\
Stranger & pupi & --- & unknown + person \\
Leader & wapi & --- & intense + person \\
Teacher & nanipi & nani & speak-mind + person \\
Student & wunipi & wuni & consume-mind + person \\
Worker & kamupi & kamu & do-move + person \\
Neighbour & kiyipi & kiyi & bond-space + person \\
Ancestor & tikipi & tiki & past-bond + person \\
\bottomrule
\end{tabular}

\subsection{Nature \& Environment}
\begin{tabular}{@{}llll@{}}
\toprule
\textbf{English} & \textbf{Sunina} & \textbf{Short} & \textbf{Logic} \\ \midrule
Sky & yiwa & --- & space + big \\
Sun & wayi & --- & intense + light \\
Moon & mawayi & --- & opposite-intense + light \\
Star & puyi & --- & unknown + light \\
Cloud & yuyiwa & yuyi & fluid + sky \\
Wind & muyiwa & --- & move + sky \\
Rain & yumu & --- & fluid + move \\
Storm & kuyumu & kuyu & force + rain \\
Snow & makuyusa & maku & opposite-force + water (soft water) \\
Air & yisa & --- & space + thing \\
River & muyusa & muyu & move + water \\
Lake & yiyusa & yiyu & space + water \\
Ocean & wayusa & wayu & intense + water \\
Mountain & wawisa & wawi & intense + earth \\
Stone / Rock & kuwisa & kuwi & force + earth (hard earth) \\
Sand & mawawisa & mawawi & opposite-intense + earth (small earth) \\
Mud & yuwisa & yuwi & fluid + earth \\
Fire & kuyi & --- & force + light \\
Smoke & mukuyi & muku & move + fire \\
Ash & tikuyi & tiku & past + fire \\
Earth & wisa & --- & heavy + thing \\
Tree & yupi & --- & nature + life \\
Flower & suyupi & suyu & good + tree \\
Seed & nuyupi & nuyu & future + tree \\
Leaf & tayupi & tayu & present + tree \\
Thunder & nakuyu & naku & speak + storm \\
Lightning & yikuyu & --- & light + storm \\
Ice / Frost & kuwayusa & kuwayu & hard + water \\
Fog / Mist & puyusa & puyu & unknown + water \\
Forest & wayupi & --- & intense + tree \\
Field / Meadow & tawisa & --- & surface + earth \\
Dawn & nuwayi & nuwa & forward + sun \\
Dusk & manuwayi & manuwa & opposite + dawn \\
Drought & mayumu & --- & opposite + rain \\
Flood & wuyumu & --- & intense-consume + rain \\
Dew & miyumu & miyu & self + rain \\
Pond & wiyusa & wiyu & self + water \\
\bottomrule
\end{tabular}

\subsection{Body \& Senses}
\begin{tabular}{@{}llll@{}}
\toprule
\textbf{English} & \textbf{Sunina} & \textbf{Short} & \textbf{Logic} \\ \midrule
Body & pisa & --- & life + thing \\
Head & nipisa & nipi & mind + body \\
Hand & kapisa & kapi & do + body \\
Foot & mupisa & mupi & move + body \\
Arm & wakapisa & wakapi & intense-do + body \\
Leg & wamupisa & wamu & intense-move + body \\
Back & tipisa & --- & behind + body \\
Mouth & napisa & --- & speak + body \\
Nose & yusipisa & yusi & fluid-sense + body \\
Tongue & nakapisa & naka & speak-do + body \\
Tooth & wukupisa & wuku & consume-force + body \\
Bone & kupisa & kupi & force + body (hard body) \\
Skin & tapisa & tapi & present/surface + body \\
Stomach & wupisa & wupi & consume + body \\
Blood & muyupisa & muyupi & move-fluid + body \\
Heart & sipisa & sipi & feel + body \\
Eye & yiwupisa & yiwupi & see + body \\
Ear & siwupisa & siwupi & sense-consume + body \\
\bottomrule
\end{tabular}

\subsection{Animals}
Compound animal names are \textbf{descriptive classifications}: they tell the listener \textit{what kind of creature} is being discussed. Multi-root modifiers ensure phonological distinctiveness between species. For everyday reference to specific species, the loanword system (\S4.16) is preferred: \textit{na} + an adapted name gives each animal a short, unique, unmistakable label. The compound form remains as a generic descriptor --- useful when encountering an unfamiliar creature or when a loanword has not yet been standardized.

\begin{tabular}{@{}llll@{}}
\toprule
\textbf{English} & \textbf{Sunina} & \textbf{Short} & \textbf{Logic} \\ \midrule
Animal / Creature & pimu & --- & life + move \\
Bird & yimupimu & yimupi & space-move + creature \\
Fish & yumupimu & yumupi & fluid-move + creature \\
Dog & sukipimu & sukipi & good-bond + creature \\
Cat & misupimu & misu & self-good + creature \\
Horse & kumupimu & kumupi & force-move + creature \\
Snake & wisapimu & wisapi & earth + creature \\
Insect & mawapimu & mawapi & opposite-intense + creature (small creature) \\
\bottomrule
\end{tabular}

\vspace{0.3em}
\textit{Note:} Short forms follow the standard collision-avoidance rule: if the first two syllables collide with an existing word, the short form extends to three syllables. For instance, \textit{yimupimu} (bird) cannot shorten to \textit{yimu} (already ``stand/rise''), so its short form is \textit{yimupi}. Specific species beyond these core entries should be borrowed as proper-noun loanwords (\textit{na} + adapted name) for everyday use.

\subsection{Food \& Drink}
\begin{tabular}{@{}llll@{}}
\toprule
\textbf{English} & \textbf{Sunina} & \textbf{Short} & \textbf{Logic} \\ \midrule
Fruit & yuwusa & --- & nature + food \\
Grain / Bread & wisawusa & wisawu & earth + food \\
Meat & pimuwusa & pimuwu & animal + food \\
Milk & piyusa & piyu & life + water \\
Salt & kuwusa & kuwu & force + food \\
Cook & wusaka & --- & food + make \\
Hungry & mawusa & mawu & no + food \\
Thirsty & mayusa & mayu & no + water \\
Delicious & suwusa & suwu & good + food \\
\bottomrule
\end{tabular}

\subsection{Core Verbs}
\textit{Note:} The verbs \textit{wuka} (eat), \textit{yuwu} (drink), and \textit{masiya} (sleep) appear under \S\,Actions \& Objects below. Bare roots such as \textit{mu}, \textit{na}, \textit{ni}, \textit{si}, and \textit{yiwu} remain valid as broad predicates in grammar examples; the lexical verbs below are narrower standardized compounds used for more specific everyday meanings.

\medskip
\noindent\textbf{Movement}
\begin{tabular}{@{}llll@{}}
\toprule
\textbf{English} & \textbf{Sunina} & \textbf{Short} & \textbf{Logic} \\ \midrule
Go / Walk & muka & --- & move + do \\
Run & kumuka & kumu & force + go \\
Come & numuka & numu & forward + go \\
Sit / Stay & mamuka & mamu & opposite + go \\
Stand / Rise & yimuka & yimu & space + go \\
Fly & yiwamuka & yiwamu & sky + go \\
Swim & yumuka & --- & fluid + go \\
Fall & wimuka & wimu & weight + go \\
\bottomrule
\end{tabular}

\medskip
\noindent\textbf{Manipulation}
\begin{tabular}{@{}llll@{}}
\toprule
\textbf{English} & \textbf{Sunina} & \textbf{Short} & \textbf{Logic} \\ \midrule
Give & nuka & --- & forward + do \\
Take & tika & --- & toward + do \\
Hold / Keep & kika & --- & bond + do \\
Open & yika & --- & space + do \\
Close & mayika & mayi & opposite + open \\
Push & nuku & --- & forward + force \\
Pull & tikuka & --- & toward + force + do \\
Build & sakika & saki & thing + bond + do \\
Break & kukika & kuki & force + bond + do \\
\bottomrule
\end{tabular}

\medskip
\noindent\textbf{Senses \& Speech}
\begin{tabular}{@{}llll@{}}
\toprule
\textbf{English} & \textbf{Sunina} & \textbf{Short} & \textbf{Logic} \\ \midrule
Look / Observe & yisika & yisi & light + feel + do \\
Listen / Attend & nasika & nasi & speak + feel + do \\
Say / Utter & namu & --- & speak + move \\
Think & nika & --- & mind + do \\
Feel & sika & --- & feel + do \\
\bottomrule
\end{tabular}

\medskip
\noindent\textbf{Life \& Time}
\begin{tabular}{@{}llll@{}}
\toprule
\textbf{English} & \textbf{Sunina} & \textbf{Short} & \textbf{Logic} \\ \midrule
Live & pika & --- & life + do \\
Die & mapika & mapi & opposite + live \\
Grow & mupika & --- & change + live \\
Begin / Start & taka & --- & present + do \\
End / Finish & mataka & mata & opposite + begin \\
\bottomrule
\end{tabular}

\medskip
\noindent\textbf{Social \& Other}
\begin{tabular}{@{}llll@{}}
\toprule
\textbf{English} & \textbf{Sunina} & \textbf{Short} & \textbf{Logic} \\ \midrule
Help & sukika & suki & good + bond + do \\
Fight & kuwaka & kuwa & force + intense + do \\
Search & nisaka & --- & mind + thing + do \\
Teach & ninaka & nina & mind + speak + do \\
Learn & sinika & sini & feel + mind + do \\
Want / Need & simu & --- & feel + move \\
Know & niki & --- & mind + bond \\
Wait & tamu & --- & time + move \\
Return & timuka & --- & past + go \\
Carry & samuka & samu & thing + go \\
\bottomrule
\end{tabular}

\medskip
\noindent\textbf{Activity \& Physical}

\noindent\resizebox{\textwidth}{!}{%
\begin{tabular}{@{}llll@{}}
\toprule
\textbf{English} & \textbf{Sunina} & \textbf{Short} & \textbf{Logic} \\ \midrule
Play & yakika & --- & joy + bond + do \\
Dance & yasika & --- & joy + feel + do \\
Sing & yanaka & --- & joy + speak + do \\
Wash / Bathe & yusuka & yusu & fluid + good + do \\
Rest & makuka & --- & opposite-force + do \\
Dig & wisaka & --- & earth + do \\
Plant (verb) & nuyuka & --- & future + nature + do \\
Cut & kutika & kuti & force + toward + do \\
Throw & nukuka & --- & forward + force + do \\
Catch & tikika & --- & toward + bond + do \\
Climb & yinumuka & yinumu & above + go \\
Jump / Leap & kunumuka & kunumu & force + forward + go \\
Turn / Rotate & simuka & --- & feel + go (shift direction) \\
Call / Summon & wanaka & wana & intense + speak + do \\
Count & nasaka & nasa & speak + thing + do \\
\bottomrule
\end{tabular}}

\subsection{Actions \& Objects}
\begin{tabular}{@{}llll@{}}
\toprule
\textbf{English} & \textbf{Sunina} & \textbf{Short} & \textbf{Logic} \\ \midrule
Eat & wuka & --- & consume + do \\
Drink & yuwu & --- & fluid + consume \\
Sleep & masiya & masi & opposite-feel + harmony \\
House & yikisa & yiki & space + rule + thing \\
Food & wusa & --- & consume + thing \\
Path / Road & muyi & --- & move + space \\
Song & yanasa & yana & joy + speak + thing \\
\bottomrule
\end{tabular}

\vspace{0.5em}
\textit{Note:} The standardized vocabulary now contains nearly 300 entries across eighteen semantic domains. New words should follow the same derivation principles and be standardized by community consensus before widespread use. An official Sunina Academy and dictionary should govern formal compound definitions and authorized short forms. Continued expansion toward 300--500 entries remains a priority; mutual intelligibility depends on speakers sharing the same compounds rather than coining ad-hoc alternatives.

\subsection{Objects \& Artefacts}
\begin{tabular}{@{}llll@{}}
\toprule
\textbf{English} & \textbf{Sunina} & \textbf{Short} & \textbf{Logic} \\ \midrule
Tool & kusa & --- & force + thing \\
Weapon & wakusa & waku & intense + tool \\
Knife & mawakusa & mawaku & small + tool \\
Rope & kimusa & kimu & bond + move + thing \\
Clothing & kipisa & --- & bond + body \\
Boat & yumusa & --- & fluid + move + thing \\
Lamp & yikusa & yiku & light + force + thing \\
Mirror & miyisa & miyi & self + light + thing \\
Gift & nukasa & --- & forward-do + thing \\
Money & kasusa & kasu & equal + good + thing \\
Vehicle & kumusa & --- & force + move + thing \\
Bed & makasa & maka & opposite-do + thing \\
Door & yikasa & --- & open + thing \\
Wall & sakisa & --- & thing + bond + thing \\
Container & kikusa & kiku & bond + force + thing \\
Key & yikakusa & yikaku & open + tool \\
\bottomrule
\end{tabular}

\medskip
\noindent\textbf{Household}

\noindent\resizebox{\textwidth}{!}{%
\begin{tabular}{@{}llll@{}}
\toprule
\textbf{English} & \textbf{Sunina} & \textbf{Short} & \textbf{Logic} \\ \midrule
Table & kasisa & kasi & equal-feel + thing (flat surface) \\
Chair / Seat & mamukasa & --- & sit + thing \\
Cup / Bowl & wukisa & wuki & consume + bond + thing \\
Pot / Cooking vessel & wuyusa & wuyu & consume + fluid + thing \\
Plate / Dish & kukasa & --- & force-do + thing (flat hard) \\
Blanket / Cover & wipisa & wipi & weight + body \\
Window & yisakisa & yisaki & light + structure \\
Roof & yinukisa & yinuki & above + structure \\
Floor & mayinukisa & --- & below + structure \\
Basket & wikisa & wiki & weight + bond + thing \\
Pillow & nipikusa & nipiku & head + tool \\
Shelf & sasisa & sasi & thing-feel + thing (display surface) \\
Broom / Brush & susikusa & susiku & clean + tool \\
Candle & yikuyisa & yikuyi & fire-light + thing \\
\bottomrule
\end{tabular}}

\subsection{Time \& Place}
\begin{tabular}{@{}llll@{}}
\toprule
\textbf{English} & \textbf{Sunina} & \textbf{Short} & \textbf{Logic} \\ \midrule
Day & yita & --- & light + time \\
Night & mayita & --- & opposite + day \\
Morning & tanu & --- & time + forward \\
Evening & tati & --- & time + past \\
Tomorrow & nuyita & nuyi & forward + day \\
Yesterday & tiyita & --- & past + day \\
Year & yutasa & yuta & nature + time + thing \\
Always & wata & --- & intense + time \\
Never & mawata & --- & opposite + always \\
City / Town & wayiki & --- & intense + space + rule \\
Now & tasi & --- & present + sense \\
Soon & nutamu & nuta & forward + wait \\
Later & puta & --- & unknown + present \\
Before (temporal) & kita & --- & bond + time \\
After (temporal) & nuki & --- & forward + bond \\
Week & manayita & manayi & seven + day \\
Month & mawayita & --- & moon + time \\
Season & yuwata & yuwa & nature + intense + time \\
Moment / Instant & sitamu & sita & sense + wait \\
\bottomrule
\end{tabular}

\subsection{Spatial \& Direction}
These compounds function as spatial postpositions (\S3.15) --- placed \textit{after} the noun they govern and before the verb.

\begin{tabular}{@{}llll@{}}
\toprule
\textbf{English} & \textbf{Sunina} & \textbf{Short} & \textbf{Logic} \\ \midrule
Above / Up & yinu & --- & space + forward \\
Below / Down & mayinu & --- & opposite + above \\
Inside / Within & niyi & --- & inner + space \\
Outside & maniyi & mani & opposite + inside \\
Near / Close to & siyi & --- & sense + space \\
Far from & masiyi & --- & opposite + near \\
Between / Among & kayi & --- & equal + space \\
Behind (spatial) & yiti & --- & space + behind \\
\bottomrule
\end{tabular}

\subsection{Quantity \& Determiners}
Quantifiers and determiners occupy the same pre-noun modifier position as adjectives (\S3.11). They combine naturally with plural \textit{-wa} and demonstratives.

\begin{tabular}{@{}llll@{}}
\toprule
\textbf{English} & \textbf{Sunina} & \textbf{Short} & \textbf{Logic} \\ \midrule
All / Every & kiwa & --- & bond + intense (totally connected) \\
Some & pusa & --- & unknown + thing \\
None & makiwa & --- & opposite + all \\
Many / Much & puwa & --- & unknown + intense (unknowably large) \\
Few & mapuwa & mapu & opposite + many \\
Enough / Sufficient & suwa & --- & good + intense \\
Part / Piece & makisa & --- & opposite-bond + thing \\
Whole / Complete & kawasa & kawa & equal + intense + thing \\
Only / Sole & sakami & saka & thing + equal + self \\
Other / Another & musa & --- & change + thing \\
Same & kaka & --- & equal + equal \\
Different & makaka & --- & opposite + same \\
\bottomrule
\end{tabular}

\vspace{0.3em}
\noindent\textit{Examples:} \\
\textit{kiwa piwa mu-ti.} \quad (All people went.) \\
\textit{pusa nisa mi yiwu-ti.} \quad (Some book I saw. $=$ I saw some book.) \\
\textit{pi musa sa yiwu-nu.} \quad (The person other thing see-will. $=$ The person will see another thing.)

\subsection{Qualities \& States}
\begin{tabular}{@{}llll@{}}
\toprule
\textbf{English} & \textbf{Sunina} & \textbf{Short} & \textbf{Logic} \\ \midrule
Bad & masu & --- & opposite + good \\
New & nusa & --- & future + thing \\
Old (age) & tisa & --- & past + thing \\
Hot & sikuyi & --- & feel + fire \\
Cold & masikuyi & masiku & opposite + hot \\
True & sunisa & suni & good + mind + thing \\
False & masunisa & masuni & opposite + true \\
Beautiful & suyisa & suyi & good + light + thing \\
Full & wasa & --- & intense + thing \\
Empty & mawasa & --- & opposite + full \\
Strong & kunu & --- & force + forward \\
Weak & makunu & --- & opposite + strong \\
Long / Tall & wamuyi & --- & intense + move + space \\
Dark & mayisa & --- & opposite light + thing \\
Wet & yuku & --- & fluid + force \\
Dry & mayuku & --- & opposite + wet \\
Clean / Pure & susi & --- & good + feel \\
Dirty & masusi & --- & opposite + clean \\
Hard (texture) & kusi & --- & force + feel \\
Soft & makusi & --- & opposite + hard \\
Heavy & wiku & --- & weight + force \\
Light (weight) & mawiku & mawi & opposite + heavy \\
Fast / Quick & muwa & --- & move + intense \\
Slow & mamuwa & --- & opposite + fast \\
Loud & nawi & --- & speak + weight \\
Quiet / Silent & manawi & --- & opposite + loud \\
Deep & puwi & --- & unknown + weight \\
Shallow & mapuwi & --- & opposite + deep \\
Wide / Broad & tawi & --- & present + weight (spread out) \\
Narrow & matawi & --- & opposite + wide \\
\bottomrule
\end{tabular}

\subsection{Emotions \& Discourse}
\begin{tabular}{@{}llll@{}}
\toprule
\textbf{English} & \textbf{Sunina} & \textbf{Short} & \textbf{Logic} \\ \midrule
Happy & yasi & --- & joy + feel \\
Sad & wisi & --- & sorrow + feel \\
Angry & wasi & --- & intense + feel \\
Afraid & musi & --- & change + feel \\
Love & suyasi & suya & good + happy \\
Hate & masuyasi & masuya & opposite + love \\
Hope & nuyasi & nuya & forward + happy \\
Peace / Harmony & yaki & --- & joy + bond \\
Pride & wanisi & wani & intense + mind + feel \\
Thank & sunaka & --- & good + speak + do \\
Sorry / Apology & wisika & --- & sorrow + feel + do \\
Farewell & yamuka & yamu & joy + go \\
\bottomrule
\end{tabular}

\subsection{Abstract \& Mental Concepts}
\begin{tabular}{@{}llll@{}}
\toprule
\textbf{English} & \textbf{Sunina} & \textbf{Short} & \textbf{Logic} \\ \midrule
Idea & ninusa & ninu & mind + forward + thing \\
Reason / Cause & kitisa & --- & bond-past + thing \\
Problem & kunisa & kuni & force + mind + thing \\
Answer / Reply & panasa & --- & question + speak + thing \\
Solution & sukunisa & suku & good + problem \\
Power & kukisa & --- & force + bond + thing \\
Freedom & makinu & --- & opposite + bond + forward \\
Wisdom & wanisa & --- & intense + mind + thing \\
Skill / Ability & sukasa & --- & good + do + thing \\
Chance / Luck & pukisa & puki & unknown + bond + thing \\
Purpose / Goal & tuni & --- & target + mind \\
Mistake / Error & manika & --- & opposite + mind + do \\
Success & katusa & katu & equal + target + thing \\
Failure & makatusa & --- & opposite + success \\
Method / Way & nimu & --- & mind + move \\
\bottomrule
\end{tabular}

\subsection{Society, Colours \& Health}

\noindent\textbf{Society}
\begin{tabular}{@{}llll@{}}
\toprule
\textbf{English} & \textbf{Sunina} & \textbf{Short} & \textbf{Logic} \\ \midrule
King / Ruler & wakipi & waki & intense + rule + person \\
Law & kisa & --- & rule + thing \\
War / Conflict & kuwasa & --- & force + intense + thing \\
Trade / Exchange & kasamu & kasa & equal + thing + move \\
Story / History & tinisa & tini & past + mind + thing \\
Dream & niyasi & niya & mind + joy + feel \\
Danger & masusa & --- & bad + thing \\
\bottomrule
\end{tabular}

\medskip
\noindent\textbf{Materials}
\begin{tabular}{@{}llll@{}}
\toprule
\textbf{English} & \textbf{Sunina} & \textbf{Short} & \textbf{Logic} \\ \midrule
Wood & yupisa & --- & nature + body + thing \\
Metal & kuwisasa & --- & force + earth + thing \\
Cloth / Fabric & kisasa & --- & bond + thing \\
Paper & nisasa & --- & mind + thing \\
Glass & yisisa & --- & light + feel + thing \\
Clay & kawisa & kawi & do + earth + thing \\
Leather & pimusasa & pimusa & body + change + thing \\
\bottomrule
\end{tabular}

\medskip
\noindent\textbf{Colours}
\begin{tabular}{@{}llll@{}}
\toprule
\textbf{English} & \textbf{Sunina} & \textbf{Short} & \textbf{Logic} \\ \midrule
Red & kuyisa & --- & fire + thing \\
Blue & yiwasa & --- & sky + thing \\
Green & yuyisa & --- & nature + light + thing \\
White & yisusa & yisu & light + good + thing \\
Black & mayisusa & mayisu & opposite + white \\
Yellow & wayisa & --- & intense + light + thing \\
Brown & wisayisa & wisayi & earth + light + thing \\
Grey & kayisa & --- & balanced + light + thing \\
Orange & kuyiwasa & kuyiwa & fire + light + intense + thing \\
\bottomrule
\end{tabular}

\medskip
\noindent\textbf{Health}
\begin{tabular}{@{}llll@{}}
\toprule
\textbf{English} & \textbf{Sunina} & \textbf{Short} & \textbf{Logic} \\ \midrule
Healthy & supisa & --- & good + body \\
Sick / Ill & masupisa & --- & opposite + healthy \\
Heal & supika & --- & good + life + do \\
Medicine & supikasa & --- & heal + thing \\
\bottomrule
\end{tabular}

\subsection{Proper Nouns \& Loanwords}
Foreign names and borrowed terms are adapted to Sunina's CV phonology:
\begin{enumerate}
    \item Map each foreign sound to its nearest Sunina consonant and vowel.
    \item Break consonant clusters by inserting the default vowel \textit{a}.
    \item Drop word-final consonants or append the default vowel \textit{a}.
\end{enumerate}
The particle \textbf{na} (word/name) precedes a proper noun to mark it as a name: \quad \textit{na Pakisi} $=$ Paris. \quad \textit{na Yuna} $=$ John.

\newpage

\section{Base-10 Mathematics}
Numbers always end in \textbf{-na}, making them instantly recognizable as mathematical terms.

\subsection{The Number Grid (0--9)}
The vowel cycle \textit{a $\rightarrow$ i $\rightarrow$ u} repeats across three consonant blocks.

\vspace{0.5em}
\begin{tabular}{@{}ll|ll|ll@{}}
\toprule
\multicolumn{2}{c|}{\textbf{S-Block (Low)}} & \multicolumn{2}{c|}{\textbf{K-Block (Mid)}} & \multicolumn{2}{c}{\textbf{M-Block (High)}} \\ \midrule
1 & sana & 4 & kana & 7 & mana \\
2 & sina & 5 & kina & 8 & mina \\
3 & suna & 6 & kuna & 9 & muna \\ \bottomrule
\end{tabular}
\quad Zero (0): \textbf{nuna} (\textit{nu} [forward] + \textit{-na} [number]: the forward-empty state --- nothing yet exists).

\vspace{0.5em}
\textit{Logic:} If you know K-block = mid and \textit{u} = third vowel, then \textit{kuna} must be 6.

\subsection{Place Values \& Operators}

\vspace{0.5em}
\begin{tabular}{@{}ll|ll@{}}
\toprule
\textbf{Multiplier} & \textbf{Word} & \textbf{Operator} & \textbf{Word} \\ \midrule
$\times 10$ & pana & $+$ (Plus) & su \\
$\times 100$ & pina & $-$ (Minus) & masu \\
$\times 1{,}000$ & puna & $\times$ (Multiply) & waka \\
& & $\div$ (Divide) & mawaka \\
& & $=$ (Equals) & ka \\ \bottomrule
\end{tabular}

\vspace{0.5em}
\noindent\textit{Examples:} \\
\textit{305:} suna pina kina. \quad (3 $\times$ 100, 5. Digit--multiplier pairs are multiplied; consecutive groups sum by juxtaposition. Zero places are omitted.) \\
\textit{Script:} $[Z\_] \; [\overline{\triangle}\cdot] \; [\bigcirc'] \; [\overline{\triangle}\cdot] \; [\square'] \; [\overline{\triangle}\cdot]$ \\[0.3em]
\textit{1 + 2 = 3:} sana su sina ka suna. \\
\textit{Script:} $[Z\cdot] \; [\overline{\triangle}\cdot] \; [Z\_] \; [Z'] \; [\overline{\triangle}\cdot] \; [\square\cdot] \; [Z\_] \; [\overline{\triangle}\cdot]$

\vspace{0.5em}
\textit{Note:} \textit{su} also means ``good'' in semantic contexts. In math, it appears strictly between number-words and is unambiguous. The \textbf{P} root (Life/Unknown) is used for placeholders (\textit{pana, pina, puna}), serving as neutral variables/positions in base math.

\section{The Featural Writing System}
Letter shapes mimic the shape of the mouth producing the sound. The consonant dictates the base geometric shape; the vowel acts as a modifier mark.

\subsection{Base Shapes (Consonants)}
\begin{itemize}
    \item \textbf{p} (Lips): $\bigcirc$ \quad \textbf{m} (Lips/Hum): $\ominus$ \quad \textbf{w} (Lip Glide): $\smile$
    \item \textbf{t} (Tongue Tip): $\triangle$ \quad \textbf{n} (Tongue/Hum): $\overline{\triangle}$ \quad \textbf{s} (Hiss): $Z$
    \item \textbf{k} (Throat): $\square$ \quad \textbf{y} (Open Throat): $\sqcup$
\end{itemize}

\subsection{Vowel Modifiers}
Each vowel occupies a fixed position within the consonant's glyph block:
\begin{itemize}
    \item \textbf{-a}: Centered dot ($\cdot$) \quad \textbf{-i}: Top vertical tick ($'$) \quad \textbf{-u}: Bottom horizontal bar ($\_$)
\end{itemize}
Each CV syllable occupies one fixed-size square block, making the script monospaced and scannable --- similar in principle to Korean Hangul.

\textit{Handwriting Guidance:} In rapid handwriting, the three vowel marks may be exaggerated for clarity: the centered dot ($\cdot$) may be drawn as a small solid circle; the top tick ($'$) as a short upward stroke extending above the glyph box; and the bottom bar ($\_$) as a small downward curve beneath it. This allographic variation preserves legibility at speed without changing the script's logic.

\textit{Font Development:} The placeholder glyphs used throughout this document (e.g., $[\bigcirc']$) are schematic representations. A dedicated Sunina font mapping each CV syllable to a proper vector glyph is planned for a future appendix.

\subsection{Complete Glyph Table}

\vspace{0.5em}
\begin{tabular}{@{}l|ccc@{}}
\toprule
 & \textbf{-a} ($\cdot$) & \textbf{-i} ($'$) & \textbf{-u} ($\_$) \\ \midrule
\textbf{p} $\bigcirc$ & $[\bigcirc\cdot]$ & $[\bigcirc']$ & $[\bigcirc\_]$ \\
\textbf{t} $\triangle$ & $[\triangle\cdot]$ & $[\triangle']$ & $[\triangle\_]$ \\
\textbf{k} $\square$ & $[\square\cdot]$ & $[\square']$ & $[\square\_]$ \\
\textbf{m} $\ominus$ & $[\ominus\cdot]$ & $[\ominus']$ & $[\ominus\_]$ \\
\textbf{n} $\overline{\triangle}$ & $[\overline{\triangle}\cdot]$ & $[\overline{\triangle}']$ & $[\overline{\triangle}\_]$ \\
\textbf{s} $Z$ & $[Z\cdot]$ & $[Z']$ & $[Z\_]$ \\
\textbf{w} $\smile$ & $[\smile\cdot]$ & $[\smile']$ & $[\smile\_]$ \\
\textbf{y} $\sqcup$ & $[\sqcup\cdot]$ & $[\sqcup']$ & $[\sqcup\_]$ \\
\bottomrule
\end{tabular}

\subsection{Orientation}
\textbf{Standard:} Horizontal left-to-right, matching the majority of the world's writing systems and all digital interfaces. \\
\textbf{Artistic/Ceremonial:} Vertical top-to-bottom (columns left-to-right), reserved for calligraphy.

\subsection{Punctuation}
\textbf{Bracket Convention:} Square brackets $[\;]$ always denote a \textbf{syllable block} (a spoken sound). Round parentheses $(\;)$ always denote a \textbf{punctuation mark} (a silent symbol). This distinction ensures that syllables and punctuation are never visually confused.

\begin{itemize}
    \item \textbf{Word boundary:} Space (gap between glyph blocks).
    \item \textbf{Sentence end:} $|$ (vertical bar).
    \item \textbf{Question:} $(\bigcirc\cdot)$ at end of sentence --- the \textit{pa} glyph in round parentheses.
    \item \textbf{Quotation:} $\|$ before and after quoted speech.
\end{itemize}

\section{Story Translation: ``The Meeting''}
\textbf{English:} \\
The person walked to the water. The person saw a fierce creature. The creature did not attack, because the creature was a friend. The person felt deep joy.

\vspace{0.5em}
\textbf{Sunina (Romanized):} \\
pi yusa nu mu-ti. $|$ \\
pi waku yiwu-ti. $|$ \\
waku ku ma ka-ti, kiti waku supi ka-ti. $|$ \\
pi wa ya si-ti. $|$

\vspace{0.5em}
\textbf{Horizontal Script (one block per syllable):} \\
\textit{S1:} $[\bigcirc'] \; [\sqcup\_] \; [Z\cdot] \; [\overline{\triangle}\_] \; [\ominus\_] \; [\triangle'] \mid$ \\
\textit{S2:} $[\bigcirc'] \; [\smile\cdot] \; [\square\_] \; [\sqcup'] \; [\smile\_] \; [\triangle'] \mid$ \\
\textit{S3:} $[\smile\cdot] \; [\square\_] \; [\square\_] \; [\ominus\cdot] \; [\square\cdot] \; [\triangle'] \; [\square'] \; [\triangle'] \; [\smile\cdot] \; [\square\_] \; [Z\_] \; [\bigcirc'] \; [\square\cdot] \; [\triangle'] \mid$ \\
\textit{S4:} $[\bigcirc'] \; [\smile\cdot] \; [\sqcup\cdot] \; [Z'] \; [\triangle'] \mid$

\end{document}